\subsection{The upper return and the necessary lower geometry}
\label{rw:finite:geometry}

The first task is to locate the possible endpoints of returning cycle
data. The zero-scale upper return determines the regular input range.
Whole-path positivity of the two lower factors and an action inequality
then exclude some combinations of losses. Use the fixed-field direction
$\ell\ge0$ chosen in Section~\ref{rw:geometry}, and put
\[
 F_\pm(X)=1\pm\bar aX+\bar BX^2,\qquad
 A_\pm(P)=\int_0^P\frac{X}{F_\pm(X)}\,\dd X
\]
on intervals where the denominator is positive throughout. Approaching
a root is allowed only from this passing side; the action is never
continued through it.

At $\epsilon=\nu=0$ the upper field, divided by its positive factor
$Y$, is
\begin{equation}\label{rw:finite:principal-upper}
 X'=\bar mY-1,\qquad Y'=X+\bar cY.
\end{equation}
Let $S''-\bar cS'-\bar mS=0$, $S(0)=0$, $S'(0)=1$, and
$I(T)=\int_0^T S(t)\,\dd t$. An arc beginning at $(p,0)$ has
$Y(t)=pS(t)-I(t)$ and $X=Y'-\bar cY$.
For $\mathcal D=\bar c^2+4\bar m\ge0$ write
$\gamma_\pm=(\bar c\pm\sqrt{\mathcal D})/2$, and define
\[
 p_c=\begin{cases}1/\gamma_+,&\mathcal D\ge0,\ \gamma_+>0,\\
 \infty,&\text{otherwise},\end{cases}
 \qquad q_c=p_c(\bar m,-\bar c).
\]

\begin{lemma}[The regular upper arc]\label{rw:finite:upper-return}
The first return of \eqref{rw:finite:principal-upper} through $Y>0$
exists exactly for $0<p<p_c$. Its exit $(-q,0)$ defines an increasing
analytic bijection $q=\mathcal U_0(p)$ from $(0,p_c)$ onto $(0,q_c)$.
If $T$ is the flight time, then
\begin{equation}\label{rw:finite:upper-derivative}
 \mathcal U_0'(p)=\frac pq e^{\bar cT},\qquad
 \mathcal U_0^{-1}(q;\bar m,\bar c)
       =\mathcal U_0(q;\bar m,-\bar c).
\end{equation}
On every compact regular coefficient and input set the positive-scale
return exists on a larger common tube and satisfies
\begin{equation}\label{rw:finite:regular-upper-error}
 \mathcal U_{\epsilon,\nu}(p)
 =\mathcal U_0(p)+O_N(\epsilon(1+|\log\epsilon|))+\nu O_N(1).
\end{equation}
Here and below $O_N$ includes selected finite ordinary input and
coefficient derivatives and Euler scale derivatives, with a finite
total order chosen before the domains. Common smaller image intervals
have bounded inverse jets in the same class.
\end{lemma}
\begin{proof}
On the first positive interval of $S$ put $W=S^2-IS'$.
Since $W'=\bar cW+S$,
\[
 W=e^{\bar cT}\int_0^T e^{-\bar ct}S(t)\,\dd t>0,
 \qquad (I/S)'=W/S^2>0.
\]
Thus the flight-time equation $p=I(T)/S(T)$ has at most one
solution. At that solution $Y'(T)=-W/S<0$, so the exit is
$(-q,0)$ with $q=W/S>0$.
For $\mathcal D<0$, $S=e^{\bar cT/2}\sin(\sqrt{-\mathcal D}T/2)
/(\sqrt{-\mathcal D}/2)$; the quotient $I/S$ tends to infinity
at the first zero of $S$. For $\mathcal D>0$ use
$S=(e^{\gamma_+T}-e^{\gamma_-T})/(\gamma_+-\gamma_-)$.
Its quotient tends to $1/\gamma_+$ if $\gamma_+>0$ and to infinity
otherwise. At $\mathcal D=0$, $S=Te^{\bar cT/2}$ gives the same
rule, including $\bar c=0$. These are exactly the stated ranges.
The determinant formula for transverse sections has normal factors
$p,q$ and divergence $\bar c$, proving the derivative identity.
Reflection gives its inverse and the range of $q$.

To estimate the positive-scale return on a common tube, split each
compact reference arc into two layers $0\le Y\le d_0$ and a middle
arc with $Y\ge d_0$.
On a layer $G=X+\bar cY$ has a fixed nonzero margin. Set
$g=\bar mY-1$, $F=\bar BX^2+\bar aX$,
$Q=\epsilon/(Y+\epsilon)$ and $r_0=Y^2/(Y+\epsilon)$.
The exact equation there is
\[
 X_Y=\frac{(1-Q)g+QF}{G+\nu r_0}.
\]
The perturbation at $\nu=0$ is $Q(F-g)/G$, whose integral is
$O(\epsilon\log(1+d_0/\epsilon))$. In the middle, use the
positive time change by $Y+\epsilon$; its field is
$(g+Q(F-g),G+\nu r_0)$. Ordinary flow dependence and the transverse
exit equation give the value estimate on larger compact tubes.
For the selected jets use the bounded-layer variation argument
\eqref{rw:app:bounded-variation}, with $E_\epsilon=\epsilon\partial_\epsilon$
and the integrable carriers $E_\epsilon Q=Q(1-Q)$,
$E_\epsilon r_0=-Qr_0$; further Euler
rows retain these carriers times bounded polynomials in $Q$.
Ordinary $\nu$ integration gives the last term of
\eqref{rw:finite:regular-upper-error}. The positive compact margin
in \eqref{rw:finite:upper-derivative} then gives the inverse jets on
fixed smaller common images.
\end{proof}

The lower maps must be expressed in the original central height to
compare the two complete paths. Their absolute size is exponentially
small, so we estimate a logarithmic amplitude after removing the action.
The common parameter domains below also contain the center contractions
of Section~5.

\begin{lemma}[Regular lower amplitudes]\label{rw:finite:regular-lower}
Fix a compact set of positive finite inputs with $F_+>0$ on $[0,p]$.
For $p$ bounded away from zero, the signed outgoing backward map
from $(p,1)$ to $X=0$ is
\begin{equation}\label{rw:finite:lower-height}
 L_+(p)=Z_+(p)-\ell H_+(p,\ell),\qquad
 Z_+(p)=a_+(p)e^{-A_+(p)/\epsilon},
\end{equation}
on a common enlarged input and coefficient domain. The logarithm of
$a_+>0$ has bounded selected jets, while
$H_+=\sqrt{\pi/2}+O_N(\sqrt\epsilon+|\ell|)$ and its required
ordinary $\ell$ derivatives are bounded. Reflection gives
$L_-(q)=Z_-(q)+\ell H_-(q,\ell)$ with the same properties.
Positive endpoints of either signed map imply a passage with $z>0$.

With the fixed weak references of Section~4 one can instead write
\begin{equation}\label{rw:finite:relative-lower}
 L_+(p)=s_++U(p),\quad L_-(q)=s_-+V_0(q),\qquad
 U=\alpha_Ue^{-A_+(p)/\epsilon},\quad
 V_0=\alpha_Ve^{-A_-(q)/\epsilon}.
\end{equation}
The positive amplitudes have bounded logarithmic jets on a small
ordinary $\nu$ box, and $s_--s_+=\rho\nu H$, $H>h_*>0$,
is independent of input. In particular $U',V_0'<0$ and their
logarithmic derivative rates have the usual action limits on every
smaller compact regular input interval.
\end{lemma}
\begin{proof}
Choose a small fixed $d_0>0$. On the regular layer $d_0\le z\le1$,
use $z$ as independent variable. Its denominator has a nonzero margin,
so $P=p+\epsilon R$ has bounded selected jets.
On the core use \eqref{rw:relative:core-field}, with
$\nu=\ell/\rho$ and $F_\epsilon=F_++\epsilon\bar m$.
Its linear part is
$(X/(\epsilon F_\epsilon)-\bar c/F_\epsilon)z
 +\ell/(\rho F_\epsilon)$.
The backward kernel $K(X,u)$ on $0\le X\le u\le P$ satisfies
\[
 |K(X,u)|\le Ce^{-c(u^2-X^2)/\epsilon},\quad
 \int_X^P|K|\frac u\epsilon\,\dd u\le C,
 \quad \int_X^P|K|\frac{\dd u}{\rho}\le C.
\]
The remainder has size
$C((X/\epsilon+1)z^2+|\ell z|/\rho)$.
Choose $d_0$ and then a fixed ordinary $\ell$ disk so that the
backward integral is a contraction. This constructs the signed map
as well. On $z=0$ the graph derivative has the sign of $\ell$, so
a trajectory with positive endpoints cannot leave and reenter $z>0$.
At $\ell=0$ the line itself is invariant.

At zero forcing,
$z(X)\le Cd_0e^{-c(P^2-X^2)/\epsilon}$, so
$\int_0^P(X/\epsilon+1)z(X)\,\dd X\le Cd_0$.
Removing the action leaves a bounded logarithm: the endpoint shift
$(A_+(P)-A_+(p))/\epsilon$ is bounded since $P-p=\epsilon R$,
the linear correction uses $F_\epsilon-F_+=O(\epsilon)$, and the
nonlinear correction is the preceding integral.
The normalized Volterra argument of
Proposition~\ref{rw:transitions:central-core} supplies its selected
jets after $X=Py$: $P$ stays in a compact positive interval, its
regular endpoint jets are bounded, and the kernel satisfies the same
weighted bound \eqref{rw:app:normalized-kernel-bound}.

Differentiate the core integral with respect to $\ell$ at zero
forcing. Up to an exponentially small layer endpoint term, its
negative outgoing value is
\[
 \int_0^P
 \exp\left[-\int_0^X \mathfrak f_z(u,z_0(u))\,\dd u\right]
 \frac{(1-z_0(X))^2}{\rho D(X,z_0(X))}\,\dd X.
\]
Near zero put $X=\rho v$. The integrand tends to $e^{-v^2/2}$;
its error is $O(\rho)$ times a fixed polynomial and a Gaussian.
Away from a fixed small central interval the positive linear rate
makes the error exponentially small. The same kernels bound higher
ordinary $\ell$ variations. Hadamard integration therefore proves
the formula for $H_+$ and its jets. Reflection proves the incoming
formula.

For \eqref{rw:finite:relative-lower}, restrict the one chosen weak
reference $z_b$ at the variable endpoint $P$. Its continuation to
$X=0$ is the same reference for all inputs. On the weighted Gaussian
core of Proposition~\ref{rw:transitions:central-core}, the difference
$w=z-z_b$ satisfies $w_X=a(X)w+b(X,w)w^2$ with
\[
 a-X/(\epsilon F_+)=O(1),\quad
 \partial_\nu a=O(1),\quad |b|\le C(X/\epsilon+1).
\]
The weighted reference estimate gives
$(X/\epsilon)|\partial_\nu z_b|=O(1)$, controlling the potentially
large coefficient term. It is this weighted bound that is needed here,
rather than bounds on the unweighted high $X$ derivatives.
The positive bounded endpoint deviation and the same normalized
Volterra estimate give $\alpha_U$; reflection gives $\alpha_V$.
The Gaussian variation above, now with $\ell=\rho\nu$, gives
opposite reference signs and gap
$\sqrt\epsilon\,\nu(\sqrt{2\pi}+O(\sqrt\epsilon))$.
Finally $-\epsilon U'/U=A_+'-\epsilon(\log\alpha_U)'$ and its
reflected identity give the positive derivative margins and logarithmic
rates, using relative rather than absolute errors.
\end{proof}

\begin{lemma}[Entry bounds and action ordering]
\label{rw:finite:action-order}
There is a fixed $\chi>0$ such that every parameter-dominant cycle
in the chosen direction has $p>\chi/2$ and $q>\chi$ at small scale.
If its incoming data stay in a compact set where $F_->0$ on $[0,q]$,
then its outgoing data stay in a compact regular lower set and
\begin{equation}\label{rw:finite:action-inequality}
 A_+(p)\le A_-(q)+C\epsilon.
\end{equation}
Consequently outgoing loss with a regular incoming lower limit is
impossible.
\end{lemma}
\begin{proof}
Take $\chi=\eta/4$, with the fixed energy threshold of Section~2,
and decrease $\eta\le1/16$ if necessary. The crossing argument in
that section gives one intersection with each half of $y=0$.
For $0\le X\le\chi$, $0\le z\le1/8$, the lower denominator
is greater than $3/4$. The cycle has $0\le\ell<3s$ and
$s<e^{-\eta^2/\epsilon}$, so variation of constants gives
\[
 z_X\le(4X/(3\epsilon)+4/3)z+4s/\rho,
 \qquad z(X)\le6s e^{2X^2/(3\epsilon)+4X/3}
                  <e^{-\eta^2/(2\epsilon)}.
\]
Thus the orbit reaches $X=\chi$. A later decrease to $\chi/2$
cannot occur while $z\le1/2$, where its horizontal velocity is
positive. In $1/2<z<1$, $z_\xi\ge\chi z/4$ as long as
$X\ge\chi/2$, and the possible decrease in $X$ is at most
$4\epsilon\log2/\chi<\chi/2$. Hence $p>\chi/2$.
For the incoming side, the reference arc with $q=\chi$ has
$F_->3/4$ and $A_-(\chi)\le\eta^2/24$. Lemma~\ref{rw:finite:regular-lower}
gives $s_{\rm ref}\ge c e^{-\eta^2/(24\epsilon)}>s$.
This arc, the intervening central segment and the segment on $z=1$
bound a flow region: the top is an entry, the central edge an exit,
and uniqueness prevents crossing the reference arc. Every $q\le\chi$
would consequently give height at least $s_{\rm ref}$, a contradiction.

Now fix incoming compact regular data, so
$s\ge ce^{-A_{\max}/\epsilon}$ for a fixed $A_{\max}$.
On a fixed outgoing strip $0<X\le R$, a turn $D=0$ is a strict
maximum of $X$. Indeed, with $u=1-z$ its equation gives
$u(1+\epsilon\bar m u)=-X(\bar a+\bar BX)$ and
$u\le2(1+R)X$. Thus $z_\xi>0$ and $D_\xi=D_z z_\xi<0$
at the turn, for small scale. Before a first rightward crossing of
$X=R$ one has $D>0$. Also $F_\epsilon-D>0$, and the forcing
term is nonnegative. On $[0,\chi]$ use
$z_X/z\ge X/(\epsilon F_\epsilon)-4/3$, and beyond $\chi$ use
$z_X/z\ge(1-\epsilon/\chi)X/(\epsilon F_\epsilon)$.
Since $z(R)\le1$, integration gives
\begin{equation}\label{rw:finite:action-budget}
 (1-\epsilon/\chi)\int_0^R\frac{X}{F_\epsilon(X)}\,\dd X
       \le-\epsilon\log s+4\chi\epsilon/3.
\end{equation}
The incoming regular arc bounds the right side by $A_{\max}+o(1)$.
Since $F_\epsilon(X)\le2+X+X^2$, the integral on the left is at
least $\frac12\log(R/2)$ for $R\ge2$. A fixed sufficiently large $R$
therefore bounds $p$.
If $F_+$ then approached zero at or before a limiting endpoint,
apply the same inequality for every $R$ before its first zero and
let $R$ approach that zero. The integral diverges for both simple
and double zeros, contradicting the bound. This argument also
excludes nearby positive polynomials with minima tending to zero.
Thus both lower data sets are compact and regular. At a closing zero,
$L_+(p)=L_-(q)$ and $\ell\ge0$ give $Z_+(p)\ge Z_-(q)$.
Taking logarithms with their bounded amplitudes proves
\eqref{rw:finite:action-inequality}.
\end{proof}

\begin{lemma}[Lower cones and simultaneous finite loss]
\label{rw:finite:lower-cones}
At a limit $(\mathbf b_*,p_*,q_*)$ of cycle data, $F_+$ and $F_-$
are nonnegative up to their respective endpoints. In particular
\[
 p_* =\infty\Longrightarrow\bar B\ge0,\ \bar a\ge-2\sqrt{\bar B},
 \qquad
 q_* =\infty\Longrightarrow\bar B\ge0,\ \bar a\le2\sqrt{\bar B}.
\]
If both finite lower limits lose positivity, then $\bar B=-\beta<0$
and their endpoints are the simple roots
\[
 p_*=(\bar a+\Delta)/(2\beta),\qquad
 q_*=(-\bar a+\Delta)/(2\beta),\qquad
 \Delta=\sqrt{\bar a^2+4\beta},\qquad p_*q_*=1/\beta.
\]
If the upper limit is regular, these roots satisfy
$\mathcal U_0(p_*)=q_*$.
\end{lemma}
\begin{proof}
At an outgoing line $X=R>0$, if $F_+(R)+\epsilon\bar m\le0$,
the horizontal velocity is negative for $0<z<1$. It prevents the
first rightward crossing needed to reach $p>R$. Reflection gives
the corresponding incoming obstruction. Passing to the limit proves
nonnegativity and the two quadratic cone conditions.
For $\bar B\ge0$, a positive root of $F_+$ requires $\bar a<0$,
whereas one of $F_-$ requires $\bar a>0$. Hence both losses force
$\bar B<0$ and the displayed roots are unique and simple. The
same barrier for the nearby exact quadratic roots prevents a cycle
endpoint from lying beyond either root. Since each limiting root
was assumed at or before its endpoint, they coincide. The upper
identity follows from the regular return convergence.
\end{proof}

A strict real upper threshold also excludes returns beyond it,
uniformly for unbounded inputs. We will use this fact at infinity.
If $\mathcal D>0$, $\gamma_+>0$ on a compact coefficient block, then for every fixed $\zeta>0$ no input
$p\ge p_c+\zeta$ returns, at a sufficiently small common scale.
To prove this without bounding $p$, choose $r>0$ between the two
eigenvalues and close to $\gamma_+$, and choose $b_0>1/r$ with
$b_0<p_c+2\zeta/3$. Then
$A=\bar m+r\bar c-r^2>0$ and $rb_0-1>0$ have fixed margins.
The limiting reference arc from $p_c+\zeta$ reaches a small
$Y=d_0$ in $X+rY-b_0>0$. On the sloping boundary of the unbounded
region $Y\ge d_0$, $X+rY-b_0\ge0$, its exact outward derivative
in the time divided by $Y+\epsilon$ is
\[
 AY+rb_0-1+Q\{F(b_0-rY)-g(Y)\}
                         +r\nu Y^2/(Y+\epsilon)>0.
\]
The error is bounded by $C_{d_0}\epsilon(1+Y)+C|\nu|Y$;
the bottom boundary points inward as well. This is a forward
invariant region. Larger inputs either reach it, by uniqueness of
the layer graphs, or escape affinely; neither returns. Reflection
gives the exit exclusion. This proof requires a strict discriminant:
it gives no exclusion at a repeated root.
\label{rw:finite:strict-barrier}

Thus two infinite endpoints require $\bar c^2+4\bar m\le0$,
because otherwise at least one reflected threshold is finite.
The remaining upper limits have $\bar m<0$, or
$\bar c=\bar m=0$. In the latter case the lower cones and primary
normalization imply $\bar B>0$. If only $p_*$ is infinite and
$q_*$ finite, the incoming lower polynomial must lose positivity,
by Lemma~\ref{rw:finite:action-order}. These restrictions classify
limits of actual cycle data. The passage constructions still use
adjacent signed coefficient domains for their common comparisons.
